% !TeX program = pdflatex
% !BIB program = biber
% =====================================================================
%  Bachelorarbeit – Hauptdatei
%  Persönliche Angaben:  meta.tex
%  Pakete & Layout:      einstellungen.tex
%  Kapitel:              kapitel/*.tex  (hier per \input eingebunden)
%  Literatur:            literatur.bib  (biblatex + biber)
% =====================================================================
\documentclass[
  12pt,                 % Schriftgröße
  a4paper,
  oneside,              % einseitig (für zweiseitigen Druck: twoside)
  open=any,
  chapterprefix=false,  % true: "Kapitel 1" über der Überschrift
  headings=normal,
  parskip=half,         % Absätze durch halben Zeilenabstand statt Einzug
  listof=totoc,         % Abbildungs-/Tabellenverzeichnis im Inhaltsverzeichnis
  bibliography=totoc,   % Literaturverzeichnis im Inhaltsverzeichnis
  numbers=noenddot,
]{scrreprt}

% =====================================================================
%  Persönliche Angaben – werden auf Titelseite, in den PDF-Metadaten und
%  in der eidesstattlichen Erklärung verwendet. Hier ändern, nicht dort.
% =====================================================================
\newcommand{\arbeitTitel}{{{TITEL}}}
\newcommand{\arbeitUntertitel}{{{UNTERTITEL}}}
\newcommand{\arbeitArt}{Bachelorarbeit}
\newcommand{\arbeitAutor}{{{AUTOR}}}
\newcommand{\arbeitMatrikel}{{{MATRIKEL}}}
\newcommand{\arbeitStudiengang}{{{STUDIENGANG}}}
\newcommand{\arbeitHochschule}{{{HOCHSCHULE}}}
\newcommand{\arbeitFakultaet}{{{FAKULTAET}}}
\newcommand{\arbeitErstpruefer}{{{ERSTPRUEFER}}}
\newcommand{\arbeitZweitpruefer}{{{ZWEITPRUEFER}}}
\newcommand{\arbeitOrt}{{{ORT}}}
\newcommand{\arbeitAbgabe}{{{ABGABE}}}

% Sprache der Arbeit (ngerman oder english) und Zitierstil (numeric, ieee,
% authoryear, apa, alphabetic) – werden in einstellungen.tex verwendet.
\newcommand{\arbeitSprache}{{{SPRACHE}}}
\newcommand{\arbeitZitierstil}{{{ZITIERSTIL}}}

% =====================================================================
%  Pakete und Layout – für pdflatex. Reihenfolge ist teilweise wichtig:
%  hyperref (fast) zuletzt, cleveref/bookmark danach.
% =====================================================================

% ---------- Zeichensatz & Sprache ----------
\usepackage[T1]{fontenc}
\usepackage[utf8]{inputenc}
\usepackage[english,\arbeitSprache]{babel}   % letzte Sprache = Hauptsprache
\usepackage{lmodern}
\usepackage{microtype}
\usepackage[autostyle=true,german=quotes]{csquotes}  % \enquote{…} → „…“

% ---------- Seitenlayout ----------
\usepackage[left=3cm,right=2.5cm,top=2.5cm,bottom=2.5cm,headsep=0.8cm,marginparwidth=2.1cm,marginparsep=0.25cm]{geometry}
\usepackage{setspace}
\onehalfspacing
\usepackage[automark,headsepline]{scrlayer-scrpage}
\clearpairofpagestyles
\ihead{\leftmark}
\ohead{\pagemark}
\pagestyle{scrheadings}
\setcounter{secnumdepth}{3}   % bis \subsubsection nummerieren
\setcounter{tocdepth}{2}      % Inhaltsverzeichnis bis \subsection

% ---------- Mathematik ----------
\usepackage{amsmath,amssymb}

% ---------- Abbildungen & Tabellen ----------
\usepackage{graphicx}
\graphicspath{{abbildungen/}}
\usepackage{float}
\usepackage[font=small,labelfont=bf,margin=0.5cm]{caption}
\usepackage{subcaption}
\usepackage{booktabs}
\usepackage{tabularx}
\usepackage{longtable}
\usepackage{multirow}
\usepackage{array}
\usepackage[table]{xcolor}
\usepackage{placeins}          % \FloatBarrier: Gleitobjekte vor dem nächsten Abschnitt ausgeben
\usepackage{pdfpages}          % \includepdf für Anhänge (Fragebögen, Datenblätter)

% ---------- Listen & Textbausteine ----------
\usepackage{enumitem}
\setlist{nosep}
\usepackage{url}
\usepackage{xspace}

% ---------- Quellcode ----------
\usepackage{listings}
\definecolor{codebg}{RGB}{247,249,248}
\definecolor{codekw}{RGB}{15,111,43}
\definecolor{codecomment}{RGB}{120,130,125}
\definecolor{codestring}{RGB}{11,110,153}
\lstset{
  basicstyle=\ttfamily\small,
  keywordstyle=\color{codekw}\bfseries,
  commentstyle=\color{codecomment}\itshape,
  stringstyle=\color{codestring},
  backgroundcolor=\color{codebg},
  numbers=left, numberstyle=\tiny\color{codecomment}, numbersep=8pt,
  frame=single, framerule=0.3pt, rulecolor=\color{gray!40},
  breaklines=true, breakatwhitespace=true,
  showstringspaces=false, tabsize=2, captionpos=b,
  inputencoding=utf8, extendedchars=true,
  literate={ä}{{\"a}}1 {ö}{{\"o}}1 {ü}{{\"u}}1 {Ä}{{\"A}}1 {Ö}{{\"O}}1 {Ü}{{\"U}}1 {ß}{{\ss}}1 {–}{{--}}1 {„}{{,,}}1 {“}{{''}}1,
}
\renewcommand{\lstlistingname}{Listing}
\renewcommand{\lstlistlistingname}{Listingverzeichnis}

% ---------- Abkürzungen ----------
% \ac{KEY} beim ersten Vorkommen ausgeschrieben, danach nur die Abkürzung.
% Liste in abkuerzungen.tex pflegen.
\usepackage[printonlyused]{acronym}

% ---------- Notizen während des Schreibens ----------
% \todo{…} am Rand, \todo[inline]{…} im Text, \missingfigure{…} als Platzhalter.
% Vor der Abgabe: Option "disable" setzen → \usepackage[disable]{todonotes}
\usepackage[colorinlistoftodos,prependcaption,textsize=footnotesize]{todonotes}

% ---------- Literatur ----------
\usepackage[
  backend=biber,
  style=\arbeitZitierstil,
  maxbibnames=99,
  maxcitenames=2,
  doi=true, isbn=false, url=true,
  date=year,
]{biblatex}
\DefineBibliographyStrings{ngerman}{andothers={et\,al\adddot}}

% ---------- Verweise (immer zuletzt) ----------
\usepackage[
  colorlinks=true,
  linkcolor=black, citecolor=black, urlcolor=blue!50!black,
  pdftitle={\arbeitTitel},
  pdfauthor={\arbeitAutor},
  pdfsubject={\arbeitArt},
  bookmarksnumbered=true,
  pdfstartview=FitH,
]{hyperref}
\usepackage{bookmark}

% ---------- Eigene Befehle ----------
\newcommand{\zB}{z.\,B.\xspace}
\newcommand{\dasheisst}{d.\,h.\xspace}   % (\dh ist in T1 bereits der Buchstabe ð)
\newcommand{\bzw}{bzw.\xspace}
\newcommand{\vgl}{vgl.\xspace}
\newcommand{\code}[1]{\texttt{#1}}

\addbibresource{literatur.bib}

\begin{document}

% ---------- Vorspann (römische Seitenzahlen) ----------
\pagenumbering{roman}
% Titelseite – Angaben kommen aus meta.tex. Layout bei Bedarf an die
% Vorgaben der Hochschule anpassen (Logo: \includegraphics[width=4cm]{logo}).
\begin{titlepage}
  \centering
  \thispagestyle{empty}

  {\large \arbeitHochschule\par}
  \vspace{0.3cm}
  {\large \arbeitFakultaet\par}
  % \vspace{1cm}\includegraphics[width=4cm]{logo}\par

  \vspace{3cm}
  {\Large\scshape \arbeitArt\par}
  \vspace{1cm}
  {\huge\bfseries \arbeitTitel\par}
  \ifx\arbeitUntertitel\empty\else
    \vspace{0.5cm}
    {\Large \arbeitUntertitel\par}
  \fi

  \vfill

  {\large vorgelegt von\par}
  \vspace{0.3cm}
  {\Large\bfseries \arbeitAutor\par}
  \vspace{0.2cm}
  {Matrikelnummer \arbeitMatrikel\par}
  {Studiengang \arbeitStudiengang\par}

  \vspace{1.5cm}
  \begin{tabular}{rl}
    Erstprüfer:in:  & \arbeitErstpruefer \\
    Zweitprüfer:in: & \arbeitZweitpruefer \\
    Abgabedatum:    & \arbeitAbgabe \\
  \end{tabular}

  \vspace{1cm}
  {\arbeitOrt, \arbeitAbgabe\par}
\end{titlepage}
\clearpage

% \input{sperrvermerk}   % nur bei Arbeiten mit vertraulichen Firmendaten

\chapter*{Kurzfassung}
\addcontentsline{toc}{chapter}{Kurzfassung}
% Kurzfassung (ca. 150–250 Wörter): Problem → Ziel → Vorgehen → Ergebnis → Bedeutung.
% Keine Zitate, keine Abkürzungen ohne Erklärung, keine Verweise auf Kapitel.
% Am Ende schreiben, wenn die Ergebnisse feststehen.

\todo[inline]{Kurzfassung zuletzt schreiben: Problem, Ziel, Vorgehen, Ergebnis, Bedeutung.}

Diese Arbeit untersucht \ldots{} Ausgangspunkt ist das Problem, dass \ldots

Ziel der Arbeit ist es, \ldots{} Dazu wurde \ldots{} entwickelt und anhand von \ldots{} evaluiert.

Die Ergebnisse zeigen, dass \ldots{} Damit leistet die Arbeit einen Beitrag zu \ldots

\medskip
\noindent\textbf{Schlagwörter:} Schlagwort 1, Schlagwort 2, Schlagwort 3


\begin{otherlanguage}{english}
\chapter*{Abstract}
\addcontentsline{toc}{chapter}{Abstract}
% English abstract – same structure as the German Kurzfassung (problem, goal,
% approach, results, contribution). Many programmes require both.

This thesis investigates \ldots{} The starting point is the problem that \ldots

The goal is to \ldots{} For this purpose, \ldots{} was developed and evaluated using \ldots

The results show that \ldots{} The thesis thereby contributes to \ldots

\medskip
\noindent\textbf{Keywords:} keyword 1, keyword 2, keyword 3

\end{otherlanguage}

\tableofcontents
\listoffigures
\listoftables
\lstlistoflistings
% Abkürzungsverzeichnis – im Text \ac{API} schreiben: beim ersten Mal erscheint
% "Application Programming Interface (API)", danach nur "API".
% Weitere Befehle: \acs{API} (nur kurz), \acl{API} (nur lang), \acp{API} (Plural).
% Nur tatsächlich verwendete Einträge werden gedruckt (Option printonlyused).
\chapter*{Abkürzungsverzeichnis}
\addcontentsline{toc}{chapter}{Abkürzungsverzeichnis}
\begin{acronym}[HTTPS]
  \acro{API}{Application Programming Interface}
  \acro{BA}{Bachelorarbeit}
  \acro{GUI}{Graphical User Interface}
  \acro{IDE}{Integrated Development Environment}
  \acro{PDF}{Portable Document Format}
  \acro{REST}{Representational State Transfer}
\end{acronym}


% ---------- Hauptteil (arabische Seitenzahlen) ----------
\clearpage
\pagenumbering{arabic}

\chapter{Einleitung}
\label{cha:einleitung}

% Die Einleitung führt vom Allgemeinen zum Konkreten: Warum ist das Thema
% relevant? Welches Problem wird gelöst? Was genau ist das Ziel? Wie ist die
% Arbeit aufgebaut? Richtwert: 5–10 % des Umfangs.

\section{Motivation}
\label{sec:motivation}

% Kontext und Relevanz des Themas – gern mit einer aktuellen Quelle belegen.
Wissenschaftliche Arbeiten werden häufig mit \LaTeX{} gesetzt, weil Gliederung,
Querverweise und Literaturverwaltung zuverlässig funktionieren~\cite{lamport1994latex}.
Gleichzeitig \ldots

\section{Problemstellung}
\label{sec:problemstellung}

% Welche konkrete Lücke oder Schwierigkeit besteht? Was passiert, wenn sie
% nicht gelöst wird?
Bisherige Ansätze \ldots{} Daraus ergibt sich das Problem, dass \ldots

\section{Zielsetzung und Forschungsfrage}
\label{sec:zielsetzung}

% Ein klar überprüfbares Ziel und eine (bis drei) Forschungsfragen. Diese
% Fragen werden in Kapitel~\ref{cha:fazit} explizit beantwortet.
Ziel dieser Arbeit ist es, \ldots{} Daraus leitet sich folgende Forschungsfrage ab:

\begin{quote}
  \textbf{FF1:} Inwiefern \ldots?
\end{quote}

Zur Beantwortung werden folgende Teilziele verfolgt:
\begin{enumerate}
  \item Analyse der Anforderungen an \ldots
  \item Konzeption und prototypische Umsetzung von \ldots
  \item Evaluation anhand von \ldots
\end{enumerate}

\section{Abgrenzung}
\label{sec:abgrenzung}

% Was wird bewusst nicht behandelt? Das schützt vor der Frage "Warum fehlt X?".
Nicht Gegenstand dieser Arbeit sind \ldots

\section{Aufbau der Arbeit}
\label{sec:aufbau}

Kapitel~\ref{cha:grundlagen} erläutert die theoretischen Grundlagen und ordnet
verwandte Arbeiten ein. Kapitel~\ref{cha:methodik} beschreibt das methodische
Vorgehen und die Anforderungen. Kapitel~\ref{cha:umsetzung} dokumentiert Konzept und
Implementierung. Kapitel~\ref{cha:evaluation} bewertet die Ergebnisse, bevor
Kapitel~\ref{cha:fazit} die Arbeit zusammenfasst und einen Ausblick gibt.

\chapter{Grundlagen und Stand der Technik}
\label{cha:grundlagen}

% Hier werden Begriffe definiert, die später gebraucht werden, und der aktuelle
% Forschungs-/Technikstand mit Quellen dargestellt. Nur das aufnehmen, was für
% Konzept, Umsetzung oder Evaluation tatsächlich benötigt wird.
% Richtwert: 20–30 % des Umfangs. Jeder Absatz mit Fremdwissen braucht ein Zitat.

\section{Begriffsdefinitionen}
\label{sec:begriffe}

% Abkürzungen mit \ac{…} einführen – beim ersten Vorkommen wird die Langform
% automatisch ausgeschrieben.
Eine \ac{API} beschreibt die Schnittstelle, über die Programme miteinander
kommunizieren~\cite{kopka2003guide}. Im Kontext dieser Arbeit wird unter \ldots{}
verstanden, dass \ldots

\section{Theoretische Grundlagen}
\label{sec:theorie}

% Modelle, Verfahren, Formeln. Gleichungen nummerieren und im Text referenzieren.
Die Laufzeit des Verfahrens lässt sich durch Gleichung~\eqref{eq:laufzeit}
abschätzen:
\begin{equation}
  T(n) = a \cdot n \log n + b
  \label{eq:laufzeit}
\end{equation}
Dabei bezeichnet $n$ die Eingabegröße und $a$, $b$ sind konstante Faktoren.

\section{Verwandte Arbeiten}
\label{sec:verwandte-arbeiten}

% Wer hat ähnliche Probleme wie gelöst? Was ist die Lücke, die diese Arbeit
% schließt? Ein Vergleich in Tabellenform ist hier oft hilfreich.
\textcite{knuth1984literate} beschreibt mit \emph{Literate Programming} einen Ansatz,
Quelltext und Dokumentation zu verbinden. Tabelle~\ref{tab:verwandte} stellt die
betrachteten Ansätze gegenüber.

\begin{table}[htbp]
  \centering
  \caption{Vergleich verwandter Ansätze}
  \label{tab:verwandte}
  \begin{tabular}{lccc}
    \toprule
    Ansatz & Lokal & Versionierung & Offen \\
    \midrule
    Webbasierter Editor & -- & \checkmark & -- \\
    Lokale IDE mit Plugin & \checkmark & (\checkmark) & \checkmark \\
    Diese Arbeit & \checkmark & \checkmark & \checkmark \\
    \bottomrule
  \end{tabular}
\end{table}

\section{Zusammenfassung}
% Kurz: Was folgt aus den Grundlagen für die eigene Arbeit?
Aus den betrachteten Arbeiten folgt, dass \ldots{} Diese Lücke adressiert das in
Kapitel~\ref{cha:umsetzung} vorgestellte Konzept.

\chapter{Methodisches Vorgehen und Anforderungen}
\label{cha:methodik}

% Wie wurde gearbeitet (Vorgehensmodell, Erhebung, Werkzeuge) und was muss die
% Lösung leisten? Anforderungen nummerieren, damit sie in der Evaluation
% wieder aufgegriffen werden können.

\section{Vorgehensmodell}
\label{sec:vorgehen}

Die Arbeit folgt einem iterativen Vorgehen aus Anforderungsanalyse, Entwurf,
Implementierung und Evaluation (\vgl Abbildung~\ref{fig:vorgehen}).

\begin{figure}[htbp]
  \centering
  % \includegraphics[width=0.8\textwidth]{vorgehen}
  \fbox{\parbox{0.78\textwidth}{\centering\vspace{1.4cm}Platzhalter: Vorgehensmodell als Grafik (abbildungen/vorgehen.png)\vspace{1.4cm}}}
  \caption{Iteratives Vorgehen der Arbeit}
  \label{fig:vorgehen}
\end{figure}

\section{Anforderungen}
\label{sec:anforderungen}

% Funktionale Anforderungen (was das System tun soll) und nicht-funktionale
% Anforderungen (wie gut: Leistung, Usability, Sicherheit …).
% Priorität: M = Muss, S = Soll, K = Kann.

\begin{table}[htbp]
  \centering
  \caption{Funktionale Anforderungen}
  \label{tab:anforderungen-funktional}
  \begin{tabularx}{\textwidth}{lXc}
    \toprule
    Nr. & Anforderung & Prio \\
    \midrule
    FA1 & Ein bestehender Projektordner lässt sich öffnen und durchsuchen. & M \\
    FA2 & Änderungen werden automatisch kompiliert und als PDF angezeigt. & M \\
    FA3 & Jede Version bleibt über eine Historie wiederherstellbar. & S \\
    FA4 & Literaturangaben lassen sich per DOI importieren. & K \\
    \bottomrule
  \end{tabularx}
\end{table}

\begin{table}[htbp]
  \centering
  \caption{Nicht-funktionale Anforderungen}
  \label{tab:anforderungen-nichtfunktional}
  \begin{tabularx}{\textwidth}{lXc}
    \toprule
    Nr. & Anforderung & Prio \\
    \midrule
    NFA1 & Die Kompilierung unterbricht den Schreibfluss nicht (asynchron). & M \\
    NFA2 & Alle Daten bleiben lokal auf dem Rechner. & M \\
    NFA3 & Die Anwendung startet in unter fünf Sekunden. & S \\
    \bottomrule
  \end{tabularx}
\end{table}

\section{Bewertungskriterien}
\label{sec:kriterien}

% Woran wird der Erfolg gemessen? Diese Kriterien strukturieren Kapitel 5.
Die Lösung wird anhand folgender Kriterien bewertet:
\begin{description}
  \item[Funktionalität] Erfüllungsgrad der Anforderungen aus Tabelle~\ref{tab:anforderungen-funktional}.
  \item[Leistung] Dauer einer Kompilierung und Reaktionszeit der Oberfläche.
  \item[Nutzbarkeit] Ergebnisse einer Nutzerbefragung (\zB System Usability Scale).
\end{description}

\chapter{Konzept und Umsetzung}
\label{cha:umsetzung}

% Erst das Konzept (Architektur, Entwurfsentscheidungen, Alternativen), dann die
% Implementierung (Technologien, zentrale Komponenten, interessante Details).
% Kein Code-Abdruck ganzer Dateien – nur aussagekräftige Ausschnitte mit Erklärung.

\section{Architektur}
\label{sec:architektur}

Abbildung~\ref{fig:architektur} zeigt die Komponenten des Systems und ihre
Schnittstellen. Die Trennung in \ldots{} und \ldots{} erlaubt \ldots

\begin{figure}[htbp]
  \centering
  \missingfigure{Architekturdiagramm einfügen (abbildungen/architektur.png)}
  \caption{Systemarchitektur}
  \label{fig:architektur}
\end{figure}

\section{Entwurfsentscheidungen}
\label{sec:entwurf}

% Für jede wichtige Entscheidung: Alternativen, Kriterien, Begründung.
Für die Datenhaltung wurden zwei Alternativen betrachtet: \ldots{} Gewählt wurde
\ldots, weil \ldots

\section{Implementierung}
\label{sec:implementierung}

Die Umsetzung erfolgte in \ldots{} Listing~\ref{lst:beispiel} zeigt die zentrale
Funktion zur \ldots

% Variante 1: Code direkt aus einer Datei einbinden (bleibt automatisch aktuell)
\lstinputlisting[language=Python, caption={Zentrale Verarbeitungsfunktion}, label={lst:beispiel}, firstline=1, lastline=20]{code/beispiel.py}

% Variante 2: Code im Text
\begin{lstlisting}[language=Python, caption={Aufruf der Funktion}, label={lst:aufruf}]
ergebnis = verarbeite(daten, schwellwert=0.5)
print(ergebnis)
\end{lstlisting}

\section{Herausforderungen}
\label{sec:herausforderungen}

% Was war schwierig und wie wurde es gelöst? Das zeigt Reflexionsfähigkeit.
Eine besondere Herausforderung stellte \ldots{} dar. \todo{Beispiel ergänzen}

\chapter{Evaluation}
\label{cha:evaluation}

% Nachweis, dass die Ziele erreicht wurden – ehrlich, mit Zahlen, mit Grenzen.
% Struktur: Aufbau der Evaluation → Ergebnisse → Diskussion → Limitationen.

\section{Evaluationsdesign}
\label{sec:eval-design}

% Was wurde wie gemessen? Testdaten, Umgebung, Teilnehmende, Ablauf.
Die Evaluation erfolgte auf einem System mit \ldots{} Als Testdaten dienten \ldots

\section{Ergebnisse}
\label{sec:ergebnisse}

Tabelle~\ref{tab:ergebnisse} fasst die Messwerte zusammen.

\begin{table}[htbp]
  \centering
  \caption{Messergebnisse (Mittelwert aus zehn Durchläufen)}
  \label{tab:ergebnisse}
  \begin{tabular}{lrr}
    \toprule
    Szenario & Dauer in s & Standardabweichung \\
    \midrule
    Kleines Projekt (10 Seiten) & 2{,}1 & 0{,}2 \\
    Mittleres Projekt (40 Seiten) & 4{,}8 & 0{,}4 \\
    Großes Projekt (120 Seiten) & 11{,}3 & 0{,}9 \\
    \bottomrule
  \end{tabular}
\end{table}

\subsection{Erfüllung der Anforderungen}
% Jede Anforderung aus Kapitel 3 abhaken oder begründen, warum nicht.
Die Muss-Anforderungen FA1 und FA2 (Tabelle~\ref{tab:anforderungen-funktional})
sind vollständig erfüllt; FA4 wurde \ldots

\section{Diskussion}
\label{sec:diskussion}

% Interpretation: Was bedeuten die Ergebnisse? Stimmen sie mit Erwartungen und
% verwandten Arbeiten (Abschnitt~\ref{sec:verwandte-arbeiten}) überein?
Die Ergebnisse zeigen, dass \ldots{} Dies deckt sich mit \textcite{knuth1984literate},
der \ldots

\section{Limitationen}
\label{sec:limitationen}

% Grenzen der Aussagekraft: kleine Stichprobe, nur eine Plattform, …
Die Aussagekraft ist begrenzt durch \ldots

\chapter{Fazit und Ausblick}
\label{cha:fazit}

% Keine neuen Inhalte mehr. Zusammenfassen, Forschungsfrage explizit beantworten,
% Beitrag benennen, offene Punkte als Ausblick.

\section{Zusammenfassung}
\label{sec:zusammenfassung}

In dieser Arbeit wurde \ldots{} Dazu wurden zunächst \ldots{} (Kapitel~\ref{cha:grundlagen}),
anschließend \ldots{} (Kapitel~\ref{cha:umsetzung}) und schließlich \ldots{}
(Kapitel~\ref{cha:evaluation}).

\section{Beantwortung der Forschungsfrage}
\label{sec:beantwortung}

Die in Abschnitt~\ref{sec:zielsetzung} gestellte Forschungsfrage \textbf{FF1} lässt sich
wie folgt beantworten: \ldots

\section{Ausblick}
\label{sec:ausblick}

% Welche Erweiterungen wären sinnvoll? Was müsste anders gemacht werden?
Offen bleibt, \ldots{} In einer Folgearbeit könnte \ldots


% ---------- Verzeichnisse & Anhang ----------
\printbibliography[heading=bibintoc,title={Literaturverzeichnis}]

\appendix
\chapter{Anhang}
\label{cha:anhang}

% Material, das den Lesefluss stören würde: Fragebögen, vollständige Listings,
% große Tabellen, Screenshots, Messprotokolle. Im Text darauf verweisen.

\section{Fragebogen}
\label{sec:anhang-fragebogen}
Der vollständige Fragebogen der Nutzerbefragung (\vgl Abschnitt~\ref{sec:eval-design}).

% Externe PDF-Seiten einbinden:
% \includepdf[pages=-, pagecommand={}]{abbildungen/fragebogen.pdf}

\section{Vollständige Messdaten}
\label{sec:anhang-messdaten}

\begin{longtable}{lrr}
  \caption{Rohdaten der Messungen}\label{tab:rohdaten} \\
  \toprule
  Durchlauf & Dauer in s & Speicher in MB \\
  \midrule
  \endfirsthead
  \toprule
  Durchlauf & Dauer in s & Speicher in MB \\
  \midrule
  \endhead
  1 & 2{,}0 & 118 \\
  2 & 2{,}2 & 121 \\
  3 & 2{,}1 & 119 \\
  \bottomrule
\end{longtable}

\section{Quellcode}
\label{sec:anhang-code}
Der vollständige Quellcode liegt im Ordner \code{code/} bei \bzw ist unter
\url{https://example.org/repository} abrufbar.


% \listoftodos   % Übersicht aller offenen \todo-Notizen (vor der Abgabe entfernen)

% Eidesstattliche Erklärung – Wortlaut ggf. durch die Vorlage der Hochschule ersetzen.
\chapter*{Eidesstattliche Erklärung}
\addcontentsline{toc}{chapter}{Eidesstattliche Erklärung}
\thispagestyle{empty}

Hiermit versichere ich, \arbeitAutor, dass ich die vorliegende \arbeitArt{} mit dem Titel
\enquote{\arbeitTitel} selbstständig und ohne unzulässige fremde Hilfe angefertigt habe.
Ich habe keine anderen als die angegebenen Quellen und Hilfsmittel benutzt. Alle Stellen,
die wörtlich oder sinngemäß aus veröffentlichten oder nicht veröffentlichten Werken
entnommen wurden, sind als solche kenntlich gemacht.

Die Arbeit wurde in gleicher oder ähnlicher Form keiner anderen Prüfungsbehörde vorgelegt
und auch nicht veröffentlicht.

% Falls KI-Werkzeuge genutzt wurden, verlangen viele Hochschulen eine Angabe:
% Bei der Erstellung dieser Arbeit wurden KI-gestützte Werkzeuge (z.\,B. …) zur
% … eingesetzt. Die inhaltliche Verantwortung liegt vollständig bei mir.

\vspace{2cm}
\noindent
\begin{tabular}{@{}p{6cm}@{\hspace{2cm}}p{6cm}@{}}
  \arbeitOrt, \arbeitAbgabe & \\[1.6cm]
  \hrulefill & \hrulefill \\
  Ort, Datum & \arbeitAutor \\
\end{tabular}


\end{document}
